\subsection{Evaluation Protocol}
\label{sec:evaluation_protocol}

We evaluate synthetic time-series data from three complementary perspectives:
generation validity, similarity to real windows, and downstream task utility.
This unified protocol is shared across the joint study even though the concrete
datasets and tasks differ across PAMAP2 and MM-Fit. Following SDForger
\cite{rousseau2025sdforger} and the TSGBench evaluation taxonomy
\cite{ang2023tsgbench}, we use feature-based and distance-based similarity
metrics to quantify fidelity, prompt-adherence metrics to quantify
controllability, and task-utility metrics to test whether generated data is
useful for downstream classification or inference. Lower similarity scores
indicate closer agreement with real data, while higher adherence and utility
scores indicate better controllability and downstream performance.

\paragraph{Generation validity.}
Before computing similarity or utility, we first check whether generated samples
can be parsed and reconstructed as valid time-series windows. For a set of
$N_{\mathrm{prompt}}$ generation prompts and $N_{\mathrm{valid}}$ valid decoded
outputs, we report the valid generation rate
\begin{equation}
    \mathrm{ValidRate}
    =
    \frac{N_{\mathrm{valid}}}{N_{\mathrm{prompt}}}.
\end{equation}
We also inspect duplicate rate and class coverage. If
$N_{\mathrm{unique}}$ is the number of unique generated windows after
protocol-specific duplicate checking, the duplicate rate is
\begin{equation}
    \mathrm{DuplicateRate}
    =
    1-\frac{N_{\mathrm{unique}}}{N_{\mathrm{valid}}}.
\end{equation}
These diagnostics explain invalid decoding, mode collapse, and missing activity
classes, but they are not substitutes for fidelity or downstream utility.

\paragraph{Similarity metrics.}
For each activity or class $a$, let
$\mathcal{R}_a=\{X_i\}_{i=1}^{n_a}$ denote real windows and
$\mathcal{G}_a=\{\tilde{X}_i\}_{i=1}^{m_a}$ denote generated windows, where
$X_i,\tilde{X}_i\in\mathbb{R}^{L\times C}$ have length $L$ and $C$ channels.
To keep computations comparable, we sample at most a fixed number of windows per
class. Feature-based metrics use the sampled real and generated sets directly;
distance-based metrics use paired real--generated windows. The final reported
score is the mean over classes.

\paragraph{Feature-based similarity.}
Marginal Distribution Difference (MDD) compares empirical marginal
distributions across time steps and channels \cite{ang2023tsgbench}. For time
step $t$ and channel $c$, let $h_{t,c}^{R}(b)$ and $h_{t,c}^{G}(b)$ be
histogram density estimates for real and generated values in bin $b$. We compute
\begin{equation}
    \mathrm{MDD}(a)
    =
    \frac{1}{LCB}
    \sum_{t=1}^{L}\sum_{c=1}^{C}\sum_{b=1}^{B}
    \left|
    h_{t,c}^{R}(b)-h_{t,c}^{G}(b)
    \right| ,
\end{equation}
where $B$ is the number of histogram bins. Autocorrelation Difference (ACD)
compares temporal dependency structure:
\begin{equation}
    \mathrm{ACD}(a)
    =
    \frac{1}{C}
    \sum_{c=1}^{C}
    \left\|
    \boldsymbol{\rho}_{c}^{R}
    -
    \boldsymbol{\rho}_{c}^{G}
    \right\|_2 ,
\end{equation}
where $\boldsymbol{\rho}_{c}$ is the autocorrelation vector for channel $c$ over
lags $0,\ldots,\tau_{\max}-1$. Skewness Difference (SD) measures the difference
in distribution asymmetry:
\begin{equation}
    \mathrm{SD}(a)
    =
    \frac{1}{C}
    \sum_{c=1}^{C}
    \left|
    \gamma_c^{R}-\gamma_c^{G}
    \right| ,
\end{equation}
where $\gamma_c=\mathbb{E}[(x_c-\mu_c)^3]/\sigma_c^3$ is channel-wise skewness.
Kurtosis Difference (KD) compares tail behavior:
\begin{equation}
    \mathrm{KD}(a)
    =
    \frac{1}{C}
    \sum_{c=1}^{C}
    \left|
    \kappa_c^{R}-\kappa_c^{G}
    \right| ,
\end{equation}
where $\kappa_c=\mathbb{E}[(x_c-\mu_c)^4]/\sigma_c^4-3$ is excess kurtosis.
These definitions follow the TSGBench/SDForger feature-based metric family
\cite{ang2023tsgbench,rousseau2025sdforger}.

\paragraph{Distance-based similarity.}
Euclidean Distance (ED) measures point-wise similarity between paired real and
generated windows:
\begin{equation}
    \mathrm{ED}(a)
    =
    \frac{1}{N_a C}
    \sum_{i=1}^{N_a}\sum_{c=1}^{C}
    \left\|
    X_{i,:,c}-\tilde{X}_{i,:,c}
    \right\|_2 .
\end{equation}
Dynamic Time Warping (DTW) measures similarity under elastic temporal alignment
\cite{sakoe1978dtw}:
\begin{equation}
    \mathrm{DTW}(a)
    =
    \frac{1}{N_a}
    \sum_{i=1}^{N_a}
    D_{\mathrm{DTW}}\!\left(X_i,\tilde{X}_i\right),
\end{equation}
where
\begin{equation}
    D_{\mathrm{DTW}}(u,v)
    =
    \min_{\pi}
    \sum_{(p,q)\in\pi}
    \left\|
    u_p-v_q
    \right\|_2
\end{equation}
over monotonic warping paths $\pi$. We also compute the SDForger shapelet
reconstruction metric, reported as SHR or SHAP-RE in the SDForger
implementation \cite{rousseau2025sdforger}. Shapelets are learned from real
windows using the SIDL reconstruction objective, then generated windows are
reconstructed using the learned shapelets. Let $S$ be the learned shapelet basis,
$\alpha_i$ the sparse coefficients, and $\mathcal{T}_{\delta_i}(S)$ the basis
shifted by offsets $\delta_i$. The reconstruction score is
\begin{equation}
    \mathrm{SHR}(a)
    =
    \frac{1}{N_a}
    \sum_{i=1}^{N_a}
    \frac{1}{2}
    \left\|
    \tilde{X}_i
    -
    \alpha_i\mathcal{T}_{\delta_i}(S)
    \right\|_2^2 .
\end{equation}
When a real-vs-real baseline is available, we additionally report normalized
ratios:
\begin{equation}
    \mathrm{ED/real}
    =
    \frac{\mathrm{ED}(\hat{x},x_{\mathrm{real}})}
    {\mathrm{ED}(x_{\mathrm{real}}^{(1)},x_{\mathrm{real}}^{(2)})},
    \qquad
    \mathrm{DTW/real}
    =
    \frac{\mathrm{DTW}(\hat{x},x_{\mathrm{real}})}
    {\mathrm{DTW}(x_{\mathrm{real}}^{(1)},x_{\mathrm{real}}^{(2)})}.
\end{equation}
These ratios are interpreted only within the same dataset, channel selection,
window length, and normalization protocol; they are not used for absolute
comparison across PAMAP2 and MM-Fit.

\paragraph{Prompt adherence.}
For descriptor-level prompt-control experiments, fidelity is not sufficient:
the generated signal must also follow the requested condition. We therefore use
Spearman rank correlation \cite{spearman1904general} between requested
descriptor values and measured descriptor values from generated signals. Given
$N$ generated samples, requested descriptor values $q_i$, and measured
descriptor values $v_i$, prompt adherence is
\begin{equation}
    \rho_s(q,v)
    =
    1
    -
    \frac{6\sum_{i=1}^{N}d_i^2}
    {N(N^2-1)},
    \qquad
    d_i=\operatorname{rank}(q_i)-\operatorname{rank}(v_i),
\end{equation}
when there are no rank ties; with ties, we compute the Pearson correlation
between the ranked variables. This metric distinguishes controllability from
raw similarity: a model may produce realistic windows while failing to follow
the requested descriptor, or vice versa.

\paragraph{Downstream utility.}
Finally, we evaluate whether synthetic data is useful for downstream activity or
exercise recognition. For PAMAP2 experiments, each window is flattened and
evaluated with a RandomForestClassifier \cite{breiman2001randomforest} under
three training settings:
\begin{equation}
    \mathcal{D}_{\mathrm{train}}
    =
    \begin{cases}
    \mathcal{D}_{\mathrm{real}}, &
    \text{real-only},\\
    \mathcal{D}_{\mathrm{syn}}, &
    \text{synthetic-only},\\
    \mathcal{D}_{\mathrm{real}}\cup\mathcal{D}_{\mathrm{syn}}, &
    \text{real-plus-synthetic}.
    \end{cases}
\end{equation}
All three settings are evaluated on held-out real samples. The real-only setting
is the supervised baseline, synthetic-only is a train-on-synthetic-test-on-real
test of standalone synthetic utility, and real-plus-synthetic measures whether
synthetic data improves data augmentation. For MM-Fit, the same utility
principle is applied to the task-specific inference protocol, including
equal-size synthetic and real-data comparisons and low-real-data fine-tuning
settings. Because the inputs, labels, and tasks differ, utility scores are
compared within each protocol rather than numerically across datasets.

We report accuracy,
\begin{equation}
    \mathrm{Accuracy}
    =
    \frac{1}{N}
    \sum_{i=1}^{N}\mathbf{1}[\hat{y}_i=y_i],
\end{equation}
and macro-F1,
\begin{equation}
    \mathrm{MacroF1}
    =
    \frac{1}{|\mathcal{Y}|}
    \sum_{y\in\mathcal{Y}}
    \frac{2\,\mathrm{Precision}_y\,\mathrm{Recall}_y}
    {\mathrm{Precision}_y+\mathrm{Recall}_y}.
\end{equation}
Confusion matrices are used where space allows to identify class-wise failure
modes. In all cases, utility evaluation tests whether the generated data
preserves enough task-relevant structure to support real downstream recognition.

% BibTeX entries needed by this subsection:
%
% @article{ang2023tsgbench,
%   title        = {{TSGBench}: Time Series Generation Benchmark},
%   author       = {Ang, Yihao and Huang, Qiang and Bao, Yifan and Tung, Anthony K. H. and Huang, Zhiyong},
%   journal      = {Proceedings of the VLDB Endowment},
%   volume       = {17},
%   number       = {3},
%   pages        = {305--318},
%   year         = {2023}
% }
%
% @inproceedings{rousseau2025sdforger,
%   title     = {Forging Time Series with Language: A Large Language Model Approach to Synthetic Data Generation},
%   author    = {Rousseau, C{\'e}cile and Boschi, Tobia and Cornacchia, Giandomenico and Salwala, Dhaval and Pascale, Alessandra and Moreno, Juan Bernabe},
%   booktitle = {Advances in Neural Information Processing Systems},
%   year      = {2025}
% }
%
% @article{sakoe1978dtw,
%   title   = {Dynamic Programming Algorithm Optimization for Spoken Word Recognition},
%   author  = {Sakoe, Hiroaki and Chiba, Seibi},
%   journal = {IEEE Transactions on Acoustics, Speech, and Signal Processing},
%   volume  = {26},
%   number  = {1},
%   pages   = {43--49},
%   year    = {1978}
% }
%
% @article{spearman1904general,
%   title   = {The Proof and Measurement of Association between Two Things},
%   author  = {Spearman, Charles},
%   journal = {The American Journal of Psychology},
%   volume  = {15},
%   number  = {1},
%   pages   = {72--101},
%   year    = {1904}
% }
%
% @article{breiman2001randomforest,
%   title   = {Random Forests},
%   author  = {Breiman, Leo},
%   journal = {Machine Learning},
%   volume  = {45},
%   number  = {1},
%   pages   = {5--32},
%   year    = {2001}
% }
